% 12-conclusion
\section{Conclusion}\label{sec:conclusion}

A bidder that buys on context has to know, inside each bid, what the reader is looking at. \sys{} answers with a profile per page, computed offline, and a slot context computed at bid time where the page reports it. Three ideas carry the design. Spend on a page follows what the page is worth, so the language model is called on fewer than a third of pages and the cascade costs tens of dollars per million. Every label carries the evidence that produced it, so a refusal can be read and checked. And relevance and suitability are different questions: the first blends the page with what is near the slot, the second takes the worst of everything on screen with the ad.

The numbers we can stand behind are the ones a test or a log can check: latency in nanoseconds to microseconds, the cost per million pages, the bound that the shortlist puts on the last tier, and buyers bidding about twice as often on pages that match. The ones we cannot yet stand behind need people or impressions: golden labels judged by hand, and the experiment of \cref{sec:protocol}, which would say whether slot context changes what ads do and how far around a slot context reaches. Both are designed. Neither has run. We expect the structure to survive what they find, because the parts they would change, the thresholds, the kernel's reach and the page weight per format, are parameters rather than structure.
